

%%%%%%%%%%%%%%%
\begin{frame}[fragile]{\textsc{Forme de base d'une requête SQL}}

La forme de base est celle d'un ``bloc'' constitué de 
trois \alert{clauses}.

\vfill

\begin{minted}[baselinestretch=.9,
               fontsize=\small,
               framesep=2mm]{sql}
select expressions
from  une_table
[where conditions]
\end{minted}

\vfill

L'interprétation est \alert{toujours} la suivante (dans l'ordre)\,:

\begin{redbullet}
   \item \sqlkw{from}\,: \alert{l'espace de recherche}\,; c'est toujours une table
     \item \sqlkw{where}\,: \alert{conditions} imposées aux nuplets
            du \sqlkw{from}
       \item \sqlkw{select}\,: construit un  \alert{nuplet-résultat} à partir de chaque nuplet
          du \sqlkw{from} qui satisfait le \sqlkw{where}
\end{redbullet}

\vfill

Dans un premier temps\,: on suppose que l'espace de recherche est une des tables de la base
\end{frame}


%%%%%%%%%%%%%%%
\begin{frame}[fragile]{Illustration}


Un premier exemple très simple.
\vfill

\begin{minted}[baselinestretch=.9,
               fontsize=\small,
               framesep=2mm]{sql}
select prénom, nom from Voyageur where idVoyageur=10
\end{minted}

\vfill

Donne une table avec un nuplet-résultat.

\vfill

\begin{tabular}{c|c}
\textbf{prénom} & \textbf{nom}\\ \hline 
Phileas & Fogg\\ \hline 
\end{tabular}

\vfill 

\begin{redbullet}
   \item L'espace de recherche est la table \texttt{Voyageur}
     \item La condition est une égalité sur un attribut (\texttt{idVoyageur=10})
     \item On construit le nuplet-résultat à partir d'un nuplet-voyageur
\end{redbullet}

\end{frame}


%%%%%%%%%%%%%%%
\begin{frame}[fragile]{D'où vient SQL\,?}

SQL ne sort pas du chapeau de quelques bricoleurs. Il a pour fondement
deux langages complémentaires.

\vfill
\begin{redbullet}
  \item \alert{L'algèbre relationnelle} (déjà vu) est un ensemble d'opérations
     applicables à des tables\,;
    \item \alert{La logique formelle} (à voir en complément)  donne à SQL son aspect \alert{déclaratif} 
    (ce que l'on veut, et pas comment on veut le calculer).
\end{redbullet}

\vfill

Nous avons choisi de baser nos explications sur l'algèbre. 

\vfill

La requête précédente correspond à l'expression\,:
$$\pi_{pr\acute{e}nom, nom}(\sigma_{idVoyageur=10} (Voyageur))$$
 
\end{frame}


%%%%%%%%%%%%%%%
\begin{frame}[fragile]{Clause \sqlkw{where} = opérateur $\sigma$ de l'algèbre}

Les conditions sont celles déjà vues pour $\sigma$


\vfill
\begin{redbullet}
  \item Comparaison ($=$, $<$, $>$, $!=$) entre un attribut et une constante
  \item Comparaison entre un attribut et un autre attribut
\end{redbullet}

\vfill

Le \sqlkw{where} est un formule propositionnelle (avec \sqlkw{and}, \sqlkw{or},
\sqlkw{not} et le parenthésage) combinant ces conditions.


\vfill
 
\begin{minted}[baselinestretch=.9,
               fontsize=\small,
               framesep=2mm]{sql}
select * from Logement 
where lieu = 'Corse' or (capacité > 50 and type='Hôtel')
\end{minted}

\end{frame}

%%%%%%%%%%%%%%%
\begin{frame}[fragile]{Clause \sqlkw{select} = opérateurs $\pi$ et $\rho$  de l'algèbre}

La projection  $\pi_{A_{1}, \cdots, A_{n}}$ correspond à 
 \texttt{\sqlkw{select} A1, .., An}
 
\vfill

Le \alert{renommage} ($\rho$ en algèbre) s'exprime avec \sqlkw{as}.

\vfill

\begin{minted}[baselinestretch=.9,
               fontsize=\small,
               framesep=2mm]{sql}
select idVoyageur as id, prénom as p,  nom as n 
from Voyageur 
where idVoyageur < 30
\end{minted}

\vfill
 \begin{tabular}{c|c|c}
\textbf{id} & \textbf{p} & \textbf{n}\\ \hline 
10 & Phileas & Fogg\\ \hline 
20 & Nicolas & Bouvier\\ \hline 
\end{tabular}
 
\vfill

\end{frame}



%%%%%%%%%%%%%%%
\begin{frame}[fragile]{Attention aux doublons}


Pour éviter le tri, SQL n'élimine pas les doublons.

\vfill


Si on le souhaite, on peut les éliminer
avec  \sqlkw{distinct}.


\vfill

\begin{center}
\begin{minipage}[b]{4cm}
\begin{minted}[baselinestretch=.9,
               fontsize=\small,
               framesep=2mm]{sql}
select type 
from Logement
\end{minted}

\begin{tabular}{c}
\textbf{type}\\ \hline 
Auberge\\ \hline 
Hôtel\\ \hline 
Gîte\\ \hline 
Hôtel\\ \hline 
\end{tabular}
\end{minipage} \hspace*{1cm} \
\begin{minipage}[b]{4cm}
\begin{minted}[baselinestretch=.9,
               fontsize=\small,
               framesep=2mm]{sql}
select distinct type 
from Logement
\end{minted}

\begin{tabular}{c}
\textbf{type}\\ \hline 
Auberge\\ \hline 
Hôtel\\ \hline 
Gîte\\ \hline 
\end{tabular}
\end{minipage}
\end{center}

\end{frame}


\begin{frame}{À retenir}

Nous savons déjà faire beaucoup de choses\,: toutes
les requêtes exprimables avec $\sigma$, $\pi$ et $\rho$

\vfill

\begin{redbullet}
  \item Retenir l'interprétation d'une requête\,: d'abord
  le   \sqlkw{from}, puis le \sqlkw{where}, et enfin le \sqlkw{select}
  \item SQL présente quelques extensions à l'algèbre (SQL permet d'utiliser des 
    fonctions par
exemple\,: voir la documentation de votre système)
  \item Le \sqlkw{distinct} est à réserver aux cas où les doublons
  sont à éliminer
\end{redbullet}

\end{frame}

